\subsection{\texorpdfstring{EUCLIDEAN DISTANCE IN \(\mathbf{R}^n\)}{EUCLIDEAN DISTANCE IN R TO THE N}}

In conformity with the general definitions the Euclidean distance between two points \(\boldsymbol{x} = (x_{i})\) and \(\boldsymbol{y} = (y_{i})\) is the number

\[
d (\boldsymbol {x}, \boldsymbol {y}) = \sqrt {\sum_ {i = 1} ^ {n} (x _ {i} - y _ {i}) ^ {2}} \geqslant 0.
\]

We recall its principal properties. The relation \(d(x, y) = o\) is equivalent to \(x = y\) . We have \(d(x, y) = d(y, x)\) ; for all scalars \(t \in \mathbf{R}\) , \(d(tx, ty) = |t|d(x, y)\) ; for all \(z \in \mathbf{R}^n\) , \(d(x + z, y + z) = d(x, y)\) ; in other words, the distance between two points is invariant under translation. The distance \(d(o, x)\) from the origin \(o\) to a point \(x\) is denoted also by \(||x||\) and is called the Euclidean norm of \(x\) (or simply the norm of \(x\) , when there is no likelihood of confusion; cf.~Chapter IX, § 3, no. 3). Then \(d(x, y) = ||y - x||\) .

For \(n = 1\) , the Euclidean distance between the points \(x, y\) of \(\mathbf{R}\) reduces to the length \(|y - x|\) of the intervals with \(x\) and \(y\) as end-points. For any \(n\) , we say that \(d(x, y) = ||y - x||\) is the length of the segments with \(x\) and \(y\) as end-points.

The Euclidean distance satisfies the triangle inequality

\[
d (\boldsymbol {x}, \boldsymbol {y}) \leqslant d (\boldsymbol {x}, \boldsymbol {z}) + d (\boldsymbol {z}, \boldsymbol {y})\tag{1}
\]

for all x, y, z in \(R^{n}\) .

We recall that the proof of (1) reduces to that of the inequality

\[
\left(\sum_ {i = 1} ^ {n} \left(x _ {i} + y _ {i}\right) ^ {2}\right) ^ {\frac {1}{2}} \leqslant \left(\sum_ {i = 1} ^ {n} x _ {i} ^ {2}\right) ^ {\frac {1}{2}} + \left(\sum_ {i = 1} ^ {n} y _ {i} ^ {2}\right) ^ {\frac {1}{2}};
\]

this in turn is equivalent to the Cauchy-Schwarz inequality

\[
\left(\sum_ {i = 1} ^ {n} x _ {i} y _ {i}\right) ^ {2} \leqslant \left(\sum_ {i = 1} ^ {n} x _ {i} ^ {2}\right) \left(\sum_ {i = 1} ^ {n} y _ {i} ^ {2}\right),
\]

which is an immediate consequence of Lagrange's identity

\[
\left(\sum_ {i = 1} ^ {n} x _ {i} ^ {2}\right) \left(\sum_ {i = 1} ^ {n} y _ {i} ^ {2}\right) - \left(\sum_ {i = 1} ^ {n} x _ {i} y _ {i}\right) ^ {2} = \frac {1}{2} \sum_ {i, j} (x _ {i} y _ {j} - x _ {j} y _ {i}) ^ {2}.
\]

This proof shows at the same time that the two sides of (1) can be equal only if z is a point of the segment with x and y as end-points.

From (1) we deduce the inequality

\[
d (\boldsymbol {x}, \boldsymbol {y}) \geqslant | d (\boldsymbol {x}, \boldsymbol {z}) - d (\boldsymbol {y}, \boldsymbol {z}) |.\tag{2}
\]

Finally, if \(\pmb{x} = (x_{i}),\pmb{y} = (y_{i})\) , we have

\[
\sup _ {1 \leqslant i \leqslant n} | x _ {i} - y _ {i} | \leqslant d (x, y) \leqslant \sqrt {n}. \sup _ {1 \leqslant i \leqslant n} | x _ {i} - y _ {i} |.\tag{3}
\]

Hence a subset A of \(\mathbf{R}^n\) is bounded (§ 1, no. 1) if and only if

\[
\sup _ {\boldsymbol {x} \in \mathbf {A}} | | \boldsymbol {x} | | <   + \infty .
\]

